\chapter{Independent Comparison of Chapters 49--59}
\label{v4-arith:w6-j:chapter}

Chapters 49--59 form eleven constructions toward the G-row
biconditional
\[
\kappa^{\mathrm{Ar}}_{\mathsf{G}}
  +(\kappa^{!})^{\mathrm{Ar}}_{\mathsf{G}}=0
\quad\Longleftrightarrow\quad
\xi(s)=\xi(1-s).
\]
Their common comparison has seven faces: source disjointness, sign
normalisation, ambient category, hypotheses, functoriality, imported
equivalences, and numerical computations.  The comparison below records only
mathematical content.  It isolates three conditional points and leaves
the rest at the strength stated in the corresponding chapters.

\section{The Eleven Local Comparisons}
\label{v4-arith:w6-j:local-comparisons}

\subsection{Chapter 49: F1.a.1, F1.a.2, and \(\mathrm{L}_{\mathrm{ACD}}\)}
\label{v4-arith:w6-j:ch49}

The Taylor--Wiles ordinary comparison uses Clozel--Harris--Taylor,
Kisin, Gee--Kisin, Emerton--Gee, and Hida on the derivation side, and
Dotto--Emerton--Gee together with Fargues--Scholze on the categorical
side.  The two source families are disjoint.  Berger--Colmez fixes the
Hodge--Tate convention, and the block-diagonal Hecke action preserves
compactly generated presentable stable \(\infty\)-categories.  The
non-ordinary extension, the non-semisimple class comparison, and the
Eisenstein extension remain the named conditional statements
\Cref{v4-arith:w5-e:conj:f1a1-full},
\Cref{v4-arith:f1a3-BH:conj:Paskunas-wild-ext}, and
\Cref{v4-arith:w5-e:conj:LACD-Eisenstein}.

The Dotto--Emerton--Gee input is used only in its fixed-central-character
\(p\geq5\) fully faithful form.  No essential-surjectivity theorem and
no global variation of the central character is imported.

\subsection{Chapter 50: F1.a.3 (A+B)}
\label{v4-arith:w6-j:ch50}

The block-additivity argument is an \(\infty\)-categorical statement in
\(\mathrm{Pr}^{L}_{st}\).  Bernstein, Bushnell--Henniart, Paskunas,
Ben-Zvi--Francis--Nadler, and Lurie supply disjoint inputs for the block
decomposition and the tensor formalism.  The co-Hochschild convention is
the one of HA~5.5.3 and is compatible with Costello--Gwilliam.

The ordinary block is controlled only after the stated deformation
hypotheses.  The wild-prime centraliser and per-block residuals are not
deduced from Paskunas alone; they are precisely the narrower conjectural
pieces named in Chapter~\ref{v4-arith:w5-f:F1a3-AB}.

\subsection{Chapter 51: F1.a.3.C}
\label{v4-arith:w6-j:ch51}

Emerton--Gee gives the ambient Noetherian formal stack.  It does not by
itself give a finite nilpotent stratification, wild-prime quiver local
models, or residual-gerbe generators.  The compact-generation theorem is
therefore conditional on the local-model input stated in
\Cref{v4-arith:w5-g:f1a3-C:thm:main}.  Under that input, recollement and
compact-generation are formal consequences of
Beilinson--Bernstein--Deligne, Gaitsgory--Rozenblyum, and
Arinkin--Gaitsgory.

The equality \(2=2\) in the \(p=3\) block is a tangent-dimension
comparison.  It is not a class-level identification of the two gerbe
directions.  The class-level statement belongs to the Paskunas--wild
comparison and is not smuggled into the compact-generation proof.

\subsection{Chapter 52: Chenevier and \(F2.c\)}
\label{v4-arith:w6-j:ch52}

The determinant-trace comparison is Chenevier's pseudocharacter
formalism.  The \(F2.c\) closure uses the elementary structure of
finite adelic test categories and does not require perfectoid
descent.  The \(AG.det\) transition is the sole non-formal point: it
requires the arithmetic \(E_{2}\)-centre gluing input isolated in
Chapter~\ref{v4-arith:w5-h:bcglue-core-and-F2c}.  Bhatt--Scholze supplies prismatic and
pro-\(\acute{e}\)tale comparison technology, but no unquoted theorem is
used as a black box for the Emerton--Gee trace comparison.

\subsection{Chapter 53: The Elementary Closure}
\label{v4-arith:w6-j:ch53}

The five elementary constructions of Chapter~\ref{v4-arith:w5-direct:closure-path}
are consequences of the earlier chapters and of the explicit formula.
They do not add a new equivalence to RH.  The only RH-equivalent
content is the Weil positivity inequality already isolated as
A-TP.  The elementary chapter is therefore a normal form for the
implication chain, not an independent proof of the Riemann Hypothesis.

\subsection{Chapter 54: Arithmetic Positselski}
\label{v4-arith:w6-j:ch54}

Pontryagin duality, Greenlees--May duality, Matlis duality, and the
Bost--Connes/Consani--Marcolli formalism close the four abelian
Positselski subitems at the stated domain.  The upgrade of the abelian
Heisenberg G-row is conditional on the precise adele-class/cobar-sheaf
identification in Consani--Marcolli.  Without that primary-source
identification the result is a conditional theorem, not a theorem in
the absolute sense.

\subsection{Chapter 55: Koszul--Petersson}
\label{v4-arith:w6-j:ch55}

The Petersson comparison is a discrete-spectrum theorem.  The
Eisenstein sector is excluded from the theorem and reappears as the
Langlands--Shahidi extension.  The Maass part is conditional on the
Ramanujan input named in Chapter~\ref{v4-arith:w10-c:chapter}.  No
continuum trace is silently folded into the discrete computation.

\subsection{Chapter 56: H-P\({}^{\mathrm{Ar}}\)}
\label{v4-arith:w6-j:ch56}

The GNS construction supplies the Hilbert-space skeleton.  It does not
construct the Hilbert--Polya operator.  The twist-existence condition is
the classical zero-line obstruction in operator form.  Thus the chapter
separates construction of the representation from positivity of the
spectral measure.

\subsection{Chapter 57: HP--Li Positivity}
\label{v4-arith:w6-j:ch57}

The Mellin--Rankin--Selberg isometry reduces the restricted Li
positivity statement to the archimedean Tate distribution.  The
archimedean Tate positivity statement is equivalent to RH and is not
proved there.  The chapter identifies the exact positivity functional
and does not replace Weil positivity by positivity on a smaller test
class.

\subsection{Chapter 58: VB\({}^{*}\)}
\label{v4-arith:w6-j:ch58}

The de Branges, Hamburger, Beurling--Malliavin, determinant, and
self-adjointness hypotheses form a chain of sufficient conditions for
the Weil inequality.  Each hypothesis is at RH strength or is explicitly
conditional.  The chapter proves the implication chain, not the
hypotheses at the far end of the chain.

\subsection{Chapter 59: \(\mathrm{R}_{\mathrm{analytic}}\)}
\label{v4-arith:w6-j:ch59}

The sharp Riemann--von Mangoldt term, the Gaussian value of the balanced
explicit formula, and the Iwasawa Artin--Rees residuals are elementary
at their stated domain.  Gaussian evaluation of the balanced explicit
formula is not Weil positivity.  The missing bridge is the full
two-sided moment asymptotic strong enough to remove the
\((\log n)^{4n}\) loss.

\section{The Three Conditional Points}
\label{v4-arith:w6-j:conditional-points}

\begin{enumerate}
\item The wild-prime compact-generation theorem of
Chapter~\ref{v4-arith:w5-g:f1a3-C} is conditional on finite nilpotent
local models and class-level residual-gerbe generators.  Emerton--Gee
alone supplies the ambient formal stack.

\item The \(AG.det\) transition of Chapter~\ref{v4-arith:w5-h:bcglue-core-and-F2c} is
conditional on arithmetic \(E_{2}\)-centre gluing.  The geometric
centres of Ben-Zvi--Francis--Nadler and Francis--Gaitsgory are not, by
themselves, the arithmetic gluing theorem.

\item The abelian Heisenberg G-row upgrade of
Chapter~\ref{v4-arith:w5-d:chapter} is conditional on the exact
Consani--Marcolli adele-class/cobar-sheaf comparison used there.
\end{enumerate}

\section{Conclusion}
\label{v4-arith:w6-j:summary}

With these three qualifications, Chapters 49--59 state precisely what
they prove.  They give independent reductions of F1, of
\(\mathrm{L}_{\mathrm{ACD}}\), of the abelian Positselski sector, and of
the discrete Koszul--Petersson identity.  The RH-equivalent assertion is
not hidden in any of these reductions; it remains the archimedean Tate
positivity, equivalently Weil positivity.
